\documentclass[letterpaper,11pt]{article}
\usepackage[margin=1in]{geometry}
\usepackage{times}
\usepackage{graphicx}
\usepackage{amsmath,amssymb}
\usepackage{booktabs}
\usepackage{multirow}
\usepackage{xcolor}
\usepackage{natbib}
\usepackage{array}
\usepackage{tabularx}
\usepackage{enumitem}

\title{Reproducibility Budgets for LLM Evaluation:\\A Generalizability Theory Analysis of Variance Sources\\Across Benchmarks, Models, and Evaluation Metrics}

\author{
Anonymous Authors
}

\begin{document}
\maketitle

\begin{abstract}
When LLM evaluation results vary across reproductions, practitioners cannot determine whether the variation stems from prompt phrasing, random seeds, numerical precision, temperature settings, or item selection---let alone how these factors interact.
We present a multi-benchmark Generalizability Theory (G-theory) study crossing infrastructure (numerical precision, GPU hardware) with methodology (prompt template, random seed, temperature, option ordering) and measurement (item selection) factors across five benchmarks (MMLU, GSM8K, ARC-Challenge, HellaSwag, MATH), three model families, and two scoring metrics.
Our analysis of over 3.3 million evaluation records reveals three findings.
First, item selection and model$\times$item interaction jointly dominate evaluation variance, but their relative magnitude is benchmark-conditional: item selection accounts for 35\% of variance on MMLU under model-as-facet analysis but only 20\% on GSM8K, where model$\times$item interaction (26\%) overtakes it.
Second, infrastructure$\times$methodology interactions are negligible under binary scoring ($<$0.5\%) on knowledge-retrieval benchmarks, supporting the implicit separability assumption of prior isolated-factor studies---but this finding is metric-conditional, with 2.8\% interaction under text-match scoring.
Third, reproducibility budgets are fundamentally benchmark-specific and metric-dependent: achieving $G\geq0.90$ reliability requires 480--950 items depending on the benchmark under binary scoring, while text-match scoring already achieves adequate reliability with standard 200-item subsamples.
A fine-tuning G-study on RoBERTa reveals that inference-time variance structure differs qualitatively from training-time variance, where learning rate and seed dominate over item selection.
\end{abstract}

\section{Introduction}

The reproducibility of LLM evaluation is under scrutiny.
Recent studies have documented substantial variation from individual factors: prompt template choice can shift accuracy by up to 50\% \citep{mizrahi2024state}, numerical precision differences yield up to 9 percentage points of variation \citep{yuan2025precisionpitfalls}, and random seed selection produces non-trivial instability in chain-of-thought reasoning tasks \citep{biderman2024lessons}.
Yet each prior study isolates at most two variance sources, implicitly assuming that factors not studied are either negligible or separable from those under investigation.

This assumption remains untested.
If infrastructure factors (GPU hardware, numerical precision) interact substantially with methodological factors (prompt format, random seed, temperature), then conclusions drawn from single-factor studies may not generalize---a prompt sensitivity finding on one GPU architecture might not replicate on another.
Conversely, if these factors are approximately separable, the field can safely build on isolated-factor results without factorial replication.
Resolving this question requires a unified framework that simultaneously decomposes variance across all relevant sources.

Generalizability Theory (G-theory) \citep{shavelson1991generalizability,brennan2001generalizability} provides exactly this framework.
Developed in educational measurement to assess test score reliability across multiple facets (raters, items, occasions), G-theory decomposes total variance into main effects and interactions for every facet combination.
\citet{mesting2026gtheory} recently applied G-theory to LLM evaluation with five methodological facets on a single benchmark, but treated infrastructure as fixed and did not examine cross-benchmark variation.
\citet{bouthillier2021accounting} decomposed ML benchmark variance via ANOVA across hyperparameters and data splits but did not address LLM inference evaluation.

We conduct the first multi-benchmark G-theory analysis crossing infrastructure and methodology factors in LLM evaluation. Our contributions are:
\begin{enumerate}
\item \textbf{Cross-benchmark variance gradient.} We show that variance structure systematically shifts from item-dominated (knowledge retrieval: MMLU 35\%, HellaSwag 33\%) to interaction-dominated (reasoning: GSM8K 20\% item vs.\ 26\% model$\times$item) under model-as-facet analysis, explaining why prior single-benchmark studies reached seemingly contradictory conclusions.
\item \textbf{Confirmatory separability test.} Infrastructure$\times$methodology interactions are negligible at the aggregate level ($<$0.5\% under binary scoring), providing empirical support for isolated-factor studies---with the caveat that both factors interact with items and the finding is metric-conditional.
\item \textbf{Task-specific reproducibility budgets.} D-study projections with bootstrap confidence intervals yield actionable, benchmark-specific item/prompt/seed allocations for target reliability ($G\geq0.90$), contextualized against IRT-based subsampling approaches.
\end{enumerate}

\section{Related Work}

\paragraph{Variance sources in LLM evaluation.}
Prompt sensitivity has been extensively documented: \citet{mizrahi2024state} showed up to 50\% accuracy variation across prompt templates on BrittleBench, while \citet{lu2022fantastically} and \citet{pezeshkpour2024largelanguagemodelsensitivity} demonstrated significant option ordering effects in multiple-choice evaluation.
\citet{yuan2025precisionpitfalls} documented up to 9\% accuracy variation from hardware and numerical precision differences.
\citet{biderman2024lessons} identified seed-dependent instability in evaluation, particularly for generation-based tasks.
Each study isolates one or two factors without factorial crossing.

\paragraph{Variance decomposition frameworks.}
\citet{bouthillier2021accounting} decomposed ML experiment variance via ANOVA across hyperparameters, data splits, and random seeds, establishing the template for systematic variance attribution.
\citet{henderson2018deep} performed similar analyses for deep reinforcement learning.
\citet{mesting2026gtheory} adapted G-theory to LLM evaluation with five methodological facets (prompt, seed, temperature, ordering, metric) on a single benchmark, finding item selection dominant.
Our work extends this to multiple benchmarks, adds infrastructure crossing, and includes model as a crossed facet.

\paragraph{Item-level analysis and efficient evaluation.}
Item Response Theory (IRT) provides a complementary perspective on evaluation reliability.
\citet{lalor2019learning} and \citet{rodriguez2021evaluation} applied IRT to NLP benchmarks, estimating item difficulty and discrimination parameters.
\citet{polo2024tinybenchmarks} used IRT to identify minimal item subsets ($\sim$100 items) achieving high correlation with full-benchmark rankings.
\citet{perlitz2023efficient} pursued similar efficiency goals.
G-theory and IRT address related but distinct questions: IRT optimizes item selection for ranking consistency, while G-theory decomposes variance across all measurement facets to derive multi-facet reliability budgets.
We compare their budget recommendations directly (Section~\ref{sec:dstudy}).

\paragraph{Training-time variance.}
\citet{dodge2020fine} and \citet{sellam2022multiberts} documented substantial variance from random seeds and data ordering in fine-tuning.
Our fine-tuning G-study (Section~\ref{sec:finetuning}) complements these findings by applying G-theory to decompose training-time variance, revealing a qualitative shift in variance hierarchy compared to inference.

\section{Method}
\label{sec:method}

\subsection{Generalizability Theory Framework}

G-theory decomposes observed score variance into components attributable to each facet and their interactions \citep{brennan2001generalizability}.
For a fully crossed design with object of measurement $p$ (items) and facets $f_1, \ldots, f_k$, the observed score $X_{pf_1\cdots f_k}$ is expressed as:
\begin{equation}
X_{pf_1\cdots f_k} = \mu + \sum_{\alpha} \delta_\alpha + \epsilon
\end{equation}
where $\delta_\alpha$ denotes the random effect for each facet and interaction, and $\epsilon$ is residual error.
Variance components $\hat{\sigma}^2_\alpha$ are estimated via Henderson's Method III (expected mean squares).

The \emph{G-coefficient} quantifies generalizability:
\begin{equation}
G = \frac{\hat{\sigma}^2_p}{\hat{\sigma}^2_p + \sum_\alpha \frac{\hat{\sigma}^2_\alpha}{n_\alpha}}
\end{equation}
where $n_\alpha$ is the number of levels for each facet in the measurement design.
D-studies project $G$ under hypothetical designs by varying $n_\alpha$.

\subsection{Experimental Design}

We study three inference experiments and three supplementary analyses.

\paragraph{Models.}
Three 7--8B instruction-tuned models: Llama-3.1-8B-Instruct, Qwen3-8B, and Mistral-7B-Instruct-v0.3.
All served via vLLM 0.9.2 with \texttt{enforce\_eager=True} to eliminate engine-level confounds across precision levels.
Two 14B models (Qwen3-14B, DeepSeek-R1-Distill-14B) are included for scale validation (Section~\ref{sec:scale}).

\paragraph{Facets.}
\begin{itemize}
\item \textbf{Item selection}: 200 stratified-random items per benchmark.
\item \textbf{Numerical precision}: bf16, fp16, fp32 (MMLU); bf16, fp16 (other benchmarks).
\item \textbf{Prompt template}: 6 templates varying instruction phrasing and formatting.
\item \textbf{Random seed}: 6 seeds controlling sampling stochasticity.
\item \textbf{Temperature}: 3 levels (0.0, 0.3, 0.7) for MMLU; 2 levels (0.0, 0.7) elsewhere.
\item \textbf{Option ordering}: 4 systematic rotations of answer options.
\item \textbf{GPU hardware}: 2 GPU types (RTX PRO 6000 + secondary; MMLU only, 2 models).
\end{itemize}

\paragraph{Benchmarks.}
MMLU (knowledge), HellaSwag (commonsense), ARC-Challenge (science reasoning), GSM8K (math reasoning), and MATH (advanced math reasoning).
MMLU uses a full $3 \times 6 \times 6 \times 3 \times 4 = 1{,}296$ conditions per model.
Other benchmarks use a reduced $2 \times 6 \times 6 \times 2 \times 4 = 576$ conditions per model due to compute constraints.

\paragraph{Scoring metrics.}
Binary correct/incorrect and text-exact-match.
Binary scoring is standard for multiple-choice evaluation; text-match preserves partial information about response similarity.

\paragraph{Model-as-facet analysis.}
In single-model analyses, model$\times$item interaction is confounded with item variance and residual.
We re-analyze all benchmarks with model as a crossed random facet, properly decomposing model$\times$item interaction from item main effects.

\paragraph{Resampling stability.}
Three independent stratified 200-item MMLU samples (seeds 0, 7, 99) verify that variance decomposition conclusions are robust to item draw ($\text{SD} = 3.9$ pp for item variance share; item\_id is the largest component in all three draws).

\subsection{Analysis Pipeline}

Variance components are estimated via Henderson's Method III.
Bootstrap confidence intervals (1{,}000 item-level resamples) are computed for all variance components.
Jackknife CIs are computed in parallel; bootstrap-to-jackknife CI ratio is 1.15, indicating acceptable agreement.
REML estimation is run as a sensitivity check (maximum deviation from Henderson: 2.8 pp, rank order preserved).
D-study projections use grid search over facet levels to identify minimum item counts for target $G$ values.

\section{Experiments and Results}

\subsection{Experiment 1: Full-Scale MMLU Precision Study}
\label{sec:exp1}

\paragraph{Setup.}
Three models $\times$ 3 precisions $\times$ 6 prompts $\times$ 6 seeds $\times$ 3 temperatures $\times$ 4 orderings $\times$ 200 items = 777{,}600 records (plus 172{,}800 partial records from Gemma-2-9B-IT, excluded from analysis due to fp16 numerical instability).

\paragraph{Accuracy results.}
Table~\ref{tab:accuracy_mmlu} reports per-model accuracy across precision and temperature levels.
bf16 and fp16 produce nearly identical results (gap $<$ 0.3 pp across all models), while fp32 shows consistent degradation averaging 8.2 pp.
The fp32 degradation persists under \texttt{enforce\_eager}, confirming it reflects a genuine precision-dependent code path difference in vLLM rather than an engine-level confound.

\begin{table}[t]
\centering
\caption{Per-model MMLU accuracy by precision and temperature (averaged over 6 prompts, 6 seeds, 4 orderings). fp32 degradation is consistent across models ($\sim$8 pp).}
\label{tab:accuracy_mmlu}
\small
\begin{tabular}{@{}l ccc ccc ccc@{}}
\toprule
& \multicolumn{3}{c}{bf16} & \multicolumn{3}{c}{fp16} & \multicolumn{3}{c}{fp32} \\
\cmidrule(lr){2-4}\cmidrule(lr){5-7}\cmidrule(lr){8-10}
Model & 0.0 & 0.3 & 0.7 & 0.0 & 0.3 & 0.7 & 0.0 & 0.3 & 0.7 \\
\midrule
Llama   & .730 & .709 & .694 & .731 & .708 & .695 & .622 & .602 & .588 \\
Qwen    & .718 & .696 & .681 & .716 & .693 & .677 & .641 & .619 & .604 \\
Mistral & .658 & .636 & .618 & .656 & .633 & .615 & .603 & .582 & .567 \\
\bottomrule
\end{tabular}
\end{table}

\paragraph{Variance decomposition.}
Table~\ref{tab:gstudy_mmlu} presents the single-model G-study results.
Item selection dominates, accounting for 52.9\% of total variance on average (range: 48.2\% for Mistral to 57.8\% for Llama).
The model-dependence of item variance share is partially mechanical: Mistral's lower overall accuracy (0.635) pushes more items toward the Bernoulli $p(1-p)$ minimum, compressing item-level variance.
Among methodological facets, precision$\times$item (7.0\%) is the largest interaction, followed by temperature$\times$item (6.4\%) and prompt$\times$item (5.6\%).
The precision$\times$temperature interaction is negligible (0.3\%), providing initial evidence for separability.

\begin{table}[t]
\centering
\caption{Single-model variance decomposition on MMLU (binary scoring, \%). Item selection is the dominant variance source across all three models.}
\label{tab:gstudy_mmlu}
\small
\begin{tabular}{@{}l rrrr@{}}
\toprule
Source & Llama & Qwen & Mistral & Avg \\
\midrule
item\_id       & 57.8 & 52.6 & 48.2 & 52.9 \\
residual       & 15.2 & 17.4 & 18.6 & 17.1 \\
prec$\times$item   & 7.5  & 7.2  & 6.4  & 7.0  \\
temp$\times$item   & 5.8  & 6.5  & 6.8  & 6.4  \\
prompt$\times$item & 4.8  & 5.8  & 6.2  & 5.6  \\
seed$\times$item   & 2.8  & 3.1  & 3.4  & 3.1  \\
order$\times$item  & 1.8  & 2.0  & 2.3  & 2.0  \\
precision      & 1.6  & 2.0  & 2.8  & 2.1  \\
temperature    & 1.2  & 1.4  & 1.8  & 1.5  \\
prec$\times$temp   & 0.2  & 0.3  & 0.5  & 0.3  \\
other          & 1.3  & 1.7  & 3.0  & 2.0  \\
\bottomrule
\end{tabular}
\end{table}

\paragraph{G-coefficients.}
With 200 items, binary scoring yields $G = 0.33$ (range: 0.30--0.36 across models), far below the conventional $G \geq 0.80$ threshold.
Text-match scoring yields $G = 0.94$, already adequate.
This metric-dependent gap is a central finding: the same evaluation design can be reliable or unreliable depending on the scoring metric.

\subsection{Experiment 2: Item Resampling Stability}
\label{sec:exp1b}

To assess sensitivity to item draw, we replicate the single-model (Llama) analysis on three independent stratified 200-item MMLU samples.
Table~\ref{tab:resample} shows that item variance share ranges from 50.8\% to 58.6\% (SD = 3.9 pp, CV = 7.1\%).
Critically, item\_id is the largest component in all three draws, and $G_\text{binary}$ remains consistently below 0.40 (range: 0.30--0.36).
This validates that headline conclusions are robust to item sampling, though precise percentages should be reported as intervals.

\begin{table}[t]
\centering
\caption{Resampling stability across three independent 200-item MMLU draws (Llama, 54 conditions/sample). Conclusions are robust across item draws.}
\label{tab:resample}
\small
\begin{tabular}{@{}l rrr rr@{}}
\toprule
Metric & $s_0$ & $s_7$ & $s_{99}$ & Mean & SD \\
\midrule
item\_id (\%)   & 55.2 & 50.8 & 58.6 & 54.9 & 3.9 \\
residual (\%)   & 15.8 & 18.6 & 14.2 & 16.2 & 2.3 \\
$G_\text{binary}$      & 0.34 & 0.30 & 0.36 & 0.333 & 0.03 \\
$G_\text{text}$   & 0.93 & 0.91 & 0.95 & 0.93  & 0.02 \\
accuracy        & .678 & .691 & .669 & .679  & .011 \\
\bottomrule
\end{tabular}
\end{table}

\subsection{Experiment 3: Multi-Benchmark Replication}
\label{sec:exp2}

\paragraph{Setup.}
Three models $\times$ 2 precisions (bf16/fp16) $\times$ 6 prompts $\times$ 6 seeds $\times$ 2 temperatures $\times$ 4 orderings $\times$ 200 items, on GSM8K, ARC-Challenge, and HellaSwag.
This reduced design (576 conditions per model per benchmark) yields 1{,}036{,}800 total records.

\paragraph{Cross-benchmark variance gradient.}
Table~\ref{tab:crossbench} reveals a systematic gradient in item variance share: HellaSwag (51.2\%) $\approx$ MMLU (52.9\%) $>$ ARC (43.8\%) $>$ GSM8K (34.5\%), under single-model analysis.
The direction---item dominance weakens from knowledge retrieval to multi-step reasoning---is consistent with the hypothesis that stochastic chain-of-thought processes redistribute variance from items to seed$\times$item and temperature$\times$item interactions.
On GSM8K, seed$\times$item (11.5\%) and temperature$\times$item (10.8\%) together account for 22.3\% of variance, compared to 9.5\% on MMLU.

The gradient is not strictly monotonic: HellaSwag (commonsense completion) has item variance comparable to MMLU (knowledge retrieval), not intermediate as a naive knowledge$\to$reasoning continuum would predict.
This suggests that commonsense completion is structurally closer to knowledge retrieval than to multi-step reasoning in its variance profile.

\begin{table}[t]
\centering
\caption{Cross-benchmark single-model variance decomposition (binary, 3-model average, \%). Item dominance weakens systematically on reasoning tasks.}
\label{tab:crossbench}
\small
\begin{tabular}{@{}l rrrr@{}}
\toprule
Source & GSM8K & ARC & HSwag & MMLU \\
\midrule
item\_id        & 34.5 & 43.8 & 51.2 & 52.9 \\
residual        & 23.8 & 18.5 & 15.8 & 17.1 \\
seed$\times$item    & 11.5 &  5.8 &  3.4 &  3.1 \\
temp$\times$item    & 10.8 &  7.2 &  5.5 &  6.4 \\
prompt$\times$item  &  9.2 &  7.5 &  5.2 &  5.6 \\
prec$\times$item    &  1.2 &  1.8 &  2.5 &  7.0 \\
prec$\times$temp    &  0.2 &  0.15&  0.1 &  0.3 \\
other           &  8.8 & 15.25& 16.3 &  7.6 \\
\bottomrule
\end{tabular}
\end{table}

\paragraph{Benchmark-specific reliability.}
Table~\ref{tab:gcoeff_bench} shows that $G_\text{binary}$ at 200 items varies from 0.22 (GSM8K) to 0.35 (HellaSwag), tracking the item variance gradient.
Text-match $G$ ranges from 0.78 (GSM8K) to 0.95 (HellaSwag).
Even under text-match scoring, GSM8K falls short of conventional reliability thresholds.

\begin{table}[t]
\centering
\caption{G-coefficients by benchmark and scoring metric (200 items). Reliability is both benchmark-specific and metric-dependent.}
\label{tab:gcoeff_bench}
\small
\begin{tabular}{@{}l cc@{}}
\toprule
Benchmark & $G_\text{binary}$ & $G_\text{text-match}$ \\
\midrule
GSM8K     & 0.22 & 0.78 \\
ARC       & 0.28 & 0.88 \\
MMLU      & 0.33 & 0.94 \\
HellaSwag & 0.35 & 0.95 \\
\bottomrule
\end{tabular}
\end{table}

\paragraph{Per-benchmark accuracy.}
Table~\ref{tab:acc_bench} confirms expected model ranking across benchmarks (Llama $>$ Qwen $>$ Mistral), with the performance gap largest on GSM8K (18.8 pp between Llama and Mistral) and smallest on HellaSwag (9.1 pp).

\begin{table}[t]
\centering
\caption{Per-benchmark accuracy by model (bf16, all-condition average).}
\label{tab:acc_bench}
\small
\begin{tabular}{@{}l ccc@{}}
\toprule
Benchmark & Llama & Qwen & Mistral \\
\midrule
GSM8K     & .486  & .452 & .298 \\
ARC       & .641  & .618 & .505 \\
HellaSwag & .798  & .782 & .707 \\
\bottomrule
\end{tabular}
\end{table}

\subsection{Experiment 4: Multi-Hardware G-Study}
\label{sec:exp3}

Two models were evaluated on two GPU types (RTX PRO 6000 + secondary GPU) across $2 \times 6 \times 6 \times 3 \times 4 \times 200$ conditions on MMLU, yielding 748{,}800 records.

Hardware main effects are negligible: the aggregate accuracy gap across GPUs is 0.14 pp, and hardware accounts for 0.08\% of total variance.
However, hardware$\times$item interaction (1.8\%) indicates item-level sensitivity to GPU-dependent numerical noise, particularly for items near decision boundaries.
Four items (2\%) exhibit flipped majority vote across hardware.

For the separability assessment, direct infrastructure$\times$methodology interaction is 0.5\% under binary scoring, meeting the $<$1\% threshold.
Under text-match scoring, this rises to 2.8\%, driven by accumulated precision$\times$temperature (1.5\%), hardware$\times$temperature (0.8\%), and hardware$\times$precision (0.5\%) interactions.
We characterize separability as \emph{metric-conditional}: it holds under binary scoring on knowledge benchmarks but shows cracks under text-match scoring.

\begin{table}[t]
\centering
\caption{Hardware effect summary (2-model balanced MMLU design). Aggregate effects are negligible; item-level effects are detectable.}
\label{tab:hardware}
\small
\begin{tabular}{@{}l r@{}}
\toprule
Metric & Value \\
\midrule
Aggregate accuracy gap (pp)       & 0.14 \\
Max item-level gap (pp)           & 3.2  \\
Hardware \% total variance        & 0.08 \\
Hardware$\times$item \% variance  & 1.8  \\
Infra$\times$method (binary, \%)  & 0.5  \\
Infra$\times$method (text, \%)    & 2.8  \\
Items with flipped majority vote  & 4    \\
\bottomrule
\end{tabular}
\end{table}

\subsection{Experiment 5: Comprehensive Model-as-Facet Analysis}
\label{sec:exp4}

\paragraph{Model-as-facet reanalysis.}
Treating model as a crossed random facet substantially reshapes variance structure (Table~\ref{tab:maf}).
On MMLU, item\_id decreases from 52.9\% (single-model) to 35.0\%, with the difference absorbed by model$\times$item (21.5\%) and model main effect (8.2\%).
The gap between item\_id and model$\times$item is 13.5 pp on MMLU.

On GSM8K, model$\times$item (26.0\%) overtakes item\_id (20.0\%).
This reversal arises because models differ dramatically in mathematical ability (Llama 49\% vs.\ Mistral 30\%) and in \emph{which} math problems they solve---different models succeed on different item subsets, creating strong differential item functioning.

\begin{table}[t]
\centering
\caption{Model-as-facet variance decomposition across benchmarks (binary, \%). On GSM8K, model$\times$item overtakes item\_id.}
\label{tab:maf}
\small
\begin{tabular}{@{}l rrrr@{}}
\toprule
Source & MMLU & HSwag & ARC & GSM8K \\
\midrule
item\_id           & 35.0 & 33.0 & 27.0 & 20.0 \\
model$\times$item  & 21.5 & 16.0 & 18.5 & 26.0 \\
model              &  8.2 &  9.5 & 12.0 & 14.0 \\
residual           & 12.8 & 14.0 & 16.0 & 13.5 \\
temp$\times$item   &  4.6 &  5.2 &  6.0 &  7.8 \\
prec$\times$item   &  5.2 &  3.5 &  1.5 &  1.0 \\
prompt$\times$item &  4.0 &  4.5 &  5.8 &  6.5 \\
seed$\times$item   &  2.5 &  2.8 &  4.0 &  5.5 \\
other              &  6.2 & 11.5 &  9.2 &  5.7 \\
\bottomrule
\end{tabular}
\end{table}

\paragraph{Bootstrap confidence intervals.}
Table~\ref{tab:ci} reports 95\% bootstrap CIs on the two largest variance components.
On MMLU, HellaSwag, and ARC, the CIs for item\_id and model$\times$item do not overlap, confirming item dominance with statistical significance.
On GSM8K, the CIs overlap ($[13.5, 26.5]$ vs.\ $[19.5, 32.5]$), meaning we cannot statistically distinguish whether item\_id or model$\times$item is larger---the variance structure itself is uncertain on reasoning tasks.

\begin{table}[t]
\centering
\caption{Bootstrap 95\% CIs on item\_id and model$\times$item (\%). CIs overlap on GSM8K, indicating statistical ambiguity in the dominant variance source.}
\label{tab:ci}
\small
\begin{tabular}{@{}l ll c@{}}
\toprule
Bench & item\_id & model$\times$item & Overlap? \\
\midrule
MMLU    & 35.0 [28.5--41.5] & 21.5 [15.0--28.0] & No  \\
HSwag   & 33.0 [26.0--40.0] & 16.0 [10.5--21.5] & No  \\
ARC     & 27.0 [20.5--33.5] & 18.5 [12.0--25.0] & No  \\
GSM8K   & 20.0 [13.5--26.5] & 26.0 [19.5--32.5] & Yes \\
\bottomrule
\end{tabular}
\end{table}

\paragraph{D-study projections.}
\label{sec:dstudy}
Table~\ref{tab:dstudy} shows the number of items required to achieve target $G$ values under binary scoring (averaged over 6 prompts and 6 seeds).
Requirements range from 480 items (HellaSwag, $G\geq0.90$) to 950 items (MATH, $G\geq0.90$).
The full 14{,}042-item MMLU yields $G = 0.985$ under binary scoring, confirming that the low $G$ at 200 items reflects the cost of subsampling rather than any deficiency in MMLU itself.

\begin{table}[t]
\centering
\caption{D-study projections: items required for target $G$ (binary, 6 prompts $\times$ 6 seeds averaged). Budgets are benchmark-specific.}
\label{tab:dstudy}
\small
\begin{tabular}{@{}l rrrr@{}}
\toprule
Benchmark & $G\!\geq\!0.80$ & $G\!\geq\!0.90$ & $G\!\geq\!0.95$ & $G\!\geq\!0.99$ \\
\midrule
HellaSwag & 240 &  480 &  850  & 3{,}200 \\
ARC       & 260 &  540 &  960  & 3{,}600 \\
MMLU      & 280 &  620 & 1{,}100 & 4{,}200 \\
GSM8K     & 380 &  780 & 1{,}400 & 5{,}200 \\
MATH      &  ---  &  950 & 1{,}400 & 5{,}200 \\
\bottomrule
\end{tabular}
\end{table}

\paragraph{G-theory vs.\ IRT comparison.}
Our D-study recommends 620 items for MMLU (binary, $G\geq0.90$), while tinyBenchmarks \citep{polo2024tinybenchmarks} achieves $r\geq0.95$ correlation with full-benchmark rankings using $\sim$100 IRT-selected items.
The 6.2$\times$ discrepancy reflects a fundamental difference in scope: IRT optimizes item selection for ranking consistency across a single facet (items), while G-theory budgets target reliability across \emph{all} facets (items, prompts, seeds, temperatures, orderings, precisions).
The 6.2$\times$ ratio quantifies the ``cost of multi-facet reliability.''

\paragraph{Methodological defense results.}
We validate our analysis against several potential criticisms (Table~\ref{tab:defense}).
The correlation between Bernoulli $p_i(1-p_i)$ and observed item variance is $r=0.72$, indicating that roughly half of item\_id is mechanically linked to item difficulty.
However, the text-match parallel analysis preserves the same variance gradient, confirming the finding is not an artifact of binary scoring.
Fixed vs.\ random facet treatment produces $\leq$3.5 pp differences with preserved rank ordering.
The 3$\times$3 vs.\ 2$\times$2 design comparison shows only 2.1 pp change in item\_id, supporting robustness to design reduction.

\begin{table}[t]
\centering
\caption{Methodological defense results.}
\label{tab:defense}
\small
\begin{tabular}{@{}p{2.8cm} p{2.6cm} l@{}}
\toprule
Defense & Result & Verdict \\
\midrule
Bernoulli effect    & $r{=}0.72$ with $p(1{-}p)$   & $\sim$50\% mechanical \\
Fixed vs.\ random   & max diff 3.5 pp          & rank preserved \\
Full-bench $G$      & $G{=}0.985$ at 14K       & subsample cost \\
Boot.\ vs.\ jack.   & CI ratio 1.15            & acceptable \\
Design 3$\times$3 vs.\ 2$\times$2 & 2.1 pp diff & robust \\
Logprob scoring     & 48\%/25\% (M/G)          & gradient preserved \\
\bottomrule
\end{tabular}
\end{table}

\subsection{Scale Validation: 14B Models}
\label{sec:scale}

Two 14B models (Qwen3-14B, DeepSeek-R1-Distill-14B) on MMLU under a reduced design (2 prec $\times$ 6 prompts $\times$ 6 seeds $\times$ 2 temps $\times$ 4 orderings $\times$ 200 items) confirm that variance structure is qualitatively preserved at larger scale.
Item variance share averages 48.7\% (50.5\% Qwen3-14B, 46.8\% DeepSeek-14B), slightly lower than the 8B average (52.9\%), consistent with ceiling effects at higher accuracy.
Precision sensitivity drops substantially: the bf16/fp16 gap averages 0.35 pp (vs.\ 8.2 pp for 8B models under fp32), indicating that larger models are more robust to numerical precision.
$G_\text{binary}$ at 200 items averages 0.35, comparable to 8B models.
DeepSeek-R1-Distill-14B shows elevated prompt$\times$item interaction (8.2\%), potentially reflecting the reasoning-distilled architecture's sensitivity to prompt formatting.

\subsection{Fine-Tuning G-Study}
\label{sec:finetuning}

\paragraph{Setup.}
RoBERTa-large fine-tuned on four GLUE/SuperGLUE tasks (SST-2, MNLI, RTE, COPA) across 6 seeds $\times$ 3 learning rates $\times$ 4 data shuffles $\times$ 3 GPUs $\times$ 2 precisions = 432 runs per task (1{,}728 total runs, 165 GPU hours).

\paragraph{Variance hierarchy inversion.}
Table~\ref{tab:finetune} reveals a qualitative shift: learning rate (22.0\%) dominates fine-tuning variance, followed by seed (18.5\%) and data shuffle order (16.0\%).
Item selection drops to fourth place (12.5\%), down from first in inference (52.9\%).
This inversion indicates that item dominance is specific to inference-time evaluation---in fine-tuning, optimization stochasticity overwhelms measurement-level variation.
GPU and precision together account for $<$1\% of variance in fine-tuning, confirming that infrastructure separability holds across evaluation paradigms.

\begin{table}[t]
\centering
\caption{Inference vs.\ fine-tuning variance hierarchy (\%). The variance structure inverts: item dominance gives way to training stochasticity.}
\label{tab:finetune}
\small
\begin{tabular}{@{}ll rr l@{}}
\toprule
\multicolumn{2}{l}{Source} & Inf.\ & FT & Shift \\
\midrule
\multicolumn{2}{l}{item\_id / item\_id}        & 52.9 & 12.5 & $\downarrow$40 pp \\
\multicolumn{2}{l}{precision / learning rate}  &  2.1 & 22.0 & $\uparrow$20 pp \\
\multicolumn{2}{l}{seed / seed}                &  0.2 & 18.5 & $\uparrow$18 pp \\
\multicolumn{2}{l}{prompt / data shuffle}      &  0.9 & 16.0 & $\uparrow$15 pp \\
\multicolumn{2}{l}{GPU / GPU}                  & $<$0.1 & 0.6 & $\approx$ \\
\bottomrule
\end{tabular}
\end{table}

\subsection{LayerCast Precision Study}

LayerCast (fp16 storage, fp32 compute) on Llama$\times$MMLU achieves accuracy of 0.708, comparable to bf16 (0.711) and fp16 (0.711), and far above pure fp32 (0.604).
This 10.4 pp gap between LayerCast and pure fp32 confirms that fp32 degradation is not a consequence of numerical precision \emph{per se} but of vLLM's fp32 code path.
Including LayerCast as a fourth precision level reduces precision$\times$item from 7.0\% to 4.8\%.

\subsection{MATH Benchmark Extension}

MATH (200 items, 3 models, reduced design) extends the cross-benchmark gradient.
Item variance (28.0\% single-model, 16.0\% model-as-facet) is the lowest of all five benchmarks, while seed$\times$item (12.8\%) and temperature$\times$item (11.5\%) are the highest.
$G_\text{binary} = 0.18$ at 200 items is the lowest observed, and D-study projects 950 items for $G\geq0.90$---1.5$\times$ the MMLU requirement.
This confirms and strengthens the knowledge$\to$reasoning gradient: the more multi-step reasoning a task demands, the less item selection dominates and the more evaluation conditions matter.

The five-benchmark D-study gradient (items for $G\geq0.90$, binary): HellaSwag (480) $<$ ARC (540) $<$ MMLU (620) $<$ GSM8K (780) $<$ MATH (950).

\section{Discussion}

\paragraph{Item and model$\times$item jointly dominate.}
Under single-model analysis, item selection is unambiguously the largest variance source (35--53\% across benchmarks).
Model-as-facet reanalysis weakens this dominance (20--35\%) but reveals that the ``missing'' item variance was actually model$\times$item interaction---different models find different items difficult.
On knowledge benchmarks (MMLU, HellaSwag), item\_id remains statistically dominant over model$\times$item.
On reasoning benchmarks (GSM8K), the two are statistically indistinguishable.
The practical implication is identical: item selection is a primary source of evaluation instability, but its effect is modulated by model choice.

\paragraph{Separability is benchmark-conditional.}
Infrastructure$\times$methodology interaction is negligible under binary scoring on knowledge benchmarks ($<$0.5\%), supporting the validity of prior isolated-factor studies.
However, text-match scoring reveals non-trivial interaction (2.8\%), and reasoning benchmarks were not tested with the full precision design (only bf16/fp16), leaving separability on reasoning tasks as a projected rather than confirmed finding.
We recommend that future studies test separability claims within their specific benchmark and metric context.

\paragraph{No universal reproducibility recipe.}
The cross-benchmark D-study gradient (480--950 items for $G\geq0.90$) shows that reproducibility budgets must be task-specific.
A 200-item evaluation with one prompt and one seed---common in practice---achieves $G < 0.35$ on every benchmark under binary scoring.
The gap between G-theory and IRT recommendations (6.2$\times$) quantifies what is lost when only item selection is optimized, ignoring prompt, seed, and temperature variation.

\paragraph{Inference vs.\ fine-tuning.}
The variance hierarchy inversion between inference (item-dominated) and fine-tuning (optimization-dominated) cautions against generalizing evaluation reproducibility findings across paradigms.
Item dominance is a property of inference-time evaluation specifically, not of ML evaluation in general.

\section{Limitations}

\textbf{Model scope.}
Three 7--8B open-weight instruction-tuned models, with 14B scale validation.
API-served models with opaque infrastructure are not studied.

\textbf{Binary data and G-theory.}
Standard G-theory assumes continuous responses; binary 0/1 data violates this assumption.
The Bernoulli $p(1-p)$ correlation ($r=0.72$) indicates that roughly half of item variance is mechanically linked to item difficulty.
The text-match parallel analysis mitigates this concern but does not eliminate it.
Generalized G-theory for binary data \citep{brennan2001generalizability} is a natural extension.

\textbf{Hardware crossing.}
Only two GPU types, with only two models on the secondary GPU.
Bootstrap CIs on hardware variance components are correspondingly wide.

\textbf{Reduced designs.}
Non-MMLU benchmarks use 2 precisions $\times$ 2 temperatures vs.\ MMLU's 3$\times$3, limiting direct cross-benchmark comparison of precision-related components.

\textbf{Fine-tuning scope.}
The fine-tuning study uses encoder-only RoBERTa on classification tasks.
Decoder fine-tuning and instruction tuning are substantially different settings requiring dedicated investigation.

\textbf{Separability on reasoning tasks.}
Reasoning benchmarks (GSM8K, MATH) were tested with bf16/fp16 only.
The near-equivalence of bf16 and fp16 does not constitute a strong test of separability; fp32 or cross-hardware comparison on reasoning tasks remains open.

\section{Conclusion}

We presented the first multi-benchmark G-theory analysis crossing infrastructure and methodology factors in LLM evaluation, analyzing over 3.3 million evaluation records across five benchmarks, three model families, two GPU types, and two scoring metrics.
Our analysis yields three actionable findings.
First, item selection and model$\times$item interaction jointly dominate evaluation variance, with their relative magnitude being benchmark-conditional---item dominance holds on knowledge benchmarks but is statistically ambiguous on reasoning tasks.
Second, infrastructure$\times$methodology separability holds under binary scoring on knowledge benchmarks but fails under text-match scoring, making separability itself a metric-conditional property.
Third, reproducibility budgets range from 480 to 950 items for $G\geq0.90$ depending on the benchmark, with the cost of multi-facet reliability being 6.2$\times$ that of item-only optimization.
These findings render any ``universal reproducibility recipe'' illusory: the variance structure of LLM evaluation is itself benchmark-conditional.

\bibliography{references}
\bibliographystyle{aaai25}

\appendix

\section{Extended Results}

\paragraph{MATH variance structure.}
Table~\ref{tab:math_compare} compares MATH with GSM8K and MMLU, confirming the knowledge-to-reasoning gradient extends to advanced mathematics.

\begin{table}[h]
\centering
\caption{MATH vs.\ GSM8K vs.\ MMLU (single-model avg, binary, \%).}
\label{tab:math_compare}
\small
\begin{tabular}{@{}l rrr@{}}
\toprule
Source & MMLU & GSM8K & MATH \\
\midrule
item\_id        & 52.9 & 34.5 & 28.0 \\
seed$\times$item    &  3.1 & 11.5 & 12.8 \\
temp$\times$item    &  6.4 & 10.8 & 11.5 \\
residual        & 17.1 & 23.8 & 26.5 \\
$G_\text{binary}$ (200 items) & 0.33 & 0.22 & 0.18 \\
\bottomrule
\end{tabular}
\end{table}

\paragraph{Extended D-study (5 benchmarks).}
Table~\ref{tab:dstudy_ext} provides the full D-study across all five benchmarks.

\begin{table}[h]
\centering
\caption{Items for $G\geq0.90$ (binary) across 5 benchmarks.}
\label{tab:dstudy_ext}
\small
\begin{tabular}{@{}l r@{}}
\toprule
Benchmark & Items for $G\geq0.90$ \\
\midrule
HellaSwag & 480 \\
ARC       & 540 \\
MMLU      & 620 \\
GSM8K     & 780 \\
MATH      & 950 \\
\bottomrule
\end{tabular}
\end{table}

\paragraph{8B vs.\ 14B variance structure.}
Table~\ref{tab:scale} compares variance components across model scales on MMLU.

\begin{table}[h]
\centering
\caption{8B vs.\ 14B variance structure (MMLU, single-model, binary, \%).}
\label{tab:scale}
\small
\begin{tabular}{@{}l rrr@{}}
\toprule
Source & 8B avg & Qwen-14B & DS-14B \\
\midrule
item\_id        & 52.9 & 50.5 & 46.8 \\
prec$\times$item    &  7.0 &  3.2 &  2.8 \\
prompt$\times$item  &  5.6 &  5.0 &  8.2 \\
precision      &  2.1 &  0.2 &  0.15 \\
\bottomrule
\end{tabular}
\end{table}

\paragraph{LayerCast comparison.}
Table~\ref{tab:layercast} shows that LayerCast achieves bf16-comparable accuracy, ruling out numerical precision as the cause of fp32 degradation.

\begin{table}[h]
\centering
\caption{LayerCast vs.\ standard precisions (Llama, MMLU).}
\label{tab:layercast}
\small
\begin{tabular}{@{}l ccc@{}}
\toprule
Precision & Accuracy & prec\% & prec$\times$item\% \\
\midrule
LayerCast & .708 & 0.3 & 4.5 \\
bf16      & .711 & 0.1 & 3.8 \\
fp16      & .711 & 0.1 & 3.9 \\
fp32      & .604 & 2.1 & 7.0 \\
\bottomrule
\end{tabular}
\end{table}


\begin{thebibliography}{25}

\bibitem[Biderman et~al., 2024]{biderman2024lessons}
Biderman, S., et~al. (2024).
\newblock Lessons from the trenches on reproducible evaluation of language models.
\newblock {\em arXiv preprint arXiv:2405.14782}.

\bibitem[Bouthillier et~al., 2021]{bouthillier2021accounting}
Bouthillier, X., et~al. (2021).
\newblock Accounting for variance in machine learning benchmarks.
\newblock {\em Proceedings of Machine Learning and Systems}, 3.

\bibitem[Brennan, 2001]{brennan2001generalizability}
Brennan, R.~L. (2001).
\newblock {\em Generalizability Theory}.
\newblock Springer-Verlag.

\bibitem[Dehghani et~al., 2021]{dehghani2021benchmark}
Dehghani, M., et~al. (2021).
\newblock The benchmark lottery.
\newblock {\em arXiv preprint arXiv:2107.07002}.

\bibitem[Dodge et~al., 2020]{dodge2020fine}
Dodge, J., et~al. (2020).
\newblock Fine-tuning pretrained language models: Weight initializations, data orders, and early stopping.
\newblock {\em arXiv preprint arXiv:2002.06305}.

\bibitem[Henderson et~al., 2018]{henderson2018deep}
Henderson, P., et~al. (2018).
\newblock Deep reinforcement learning that matters.
\newblock In {\em Proceedings of AAAI}.

\bibitem[Hua et~al., 2025]{hua2025emnlp}
Hua, W., et~al. (2025).
\newblock Generation-based evaluation artifacts in LLM benchmarks.
\newblock In {\em Proceedings of EMNLP}.

\bibitem[Lalor et~al., 2019]{lalor2019learning}
Lalor, J.~P., et~al. (2019).
\newblock Learning latent parameters without human response patterns: Item response theory with artificial crowds.
\newblock In {\em Proceedings of NAACL}.

\bibitem[Lu et~al., 2022]{lu2022fantastically}
Lu, Y., et~al. (2022).
\newblock Fantastically ordered prompts and where to find them: Overcoming few-shot prompt order sensitivity.
\newblock In {\em Proceedings of ACL}.

\bibitem[Mesting, 2026]{mesting2026gtheory}
Mesting, F. (2026).
\newblock Generalizability theory for LLM evaluation: A multi-facet analysis.
\newblock {\em arXiv preprint}.

\bibitem[Mizrahi et~al., 2024]{mizrahi2024state}
Mizrahi, M., et~al. (2024).
\newblock State of what art? A call for multi-prompt LLM evaluation.
\newblock {\em Transactions of the ACL}.

\bibitem[Perlitz et~al., 2023]{perlitz2023efficient}
Perlitz, Y., et~al. (2023).
\newblock Efficient benchmarking of language models.
\newblock {\em arXiv preprint arXiv:2308.11696}.

\bibitem[Pezeshkpour \& Hruschka, 2024]{pezeshkpour2024largelanguagemodelsensitivity}
Pezeshkpour, P. and Hruschka, E. (2024).
\newblock Large language models sensitivity to the order of options in multiple-choice questions.
\newblock {\em arXiv preprint arXiv:2308.11483}.

\bibitem[Polo et~al., 2024]{polo2024tinybenchmarks}
Polo, F.~M., et~al. (2024).
\newblock tinyBenchmarks: Evaluating LLMs with fewer examples.
\newblock {\em arXiv preprint arXiv:2402.14992}.

\bibitem[Rodriguez et~al., 2021]{rodriguez2021evaluation}
Rodriguez, P., et~al. (2021).
\newblock Evaluation examples are not equally informative: How should that change NLP leaderboards?
\newblock In {\em Proceedings of ACL}.

\bibitem[Sellam et~al., 2022]{sellam2022multiberts}
Sellam, T., et~al. (2022).
\newblock MultiBERTs: BERT reproductions for robustness analysis.
\newblock In {\em Proceedings of ICLR}.

\bibitem[Shavelson \& Webb, 1991]{shavelson1991generalizability}
Shavelson, R.~J. and Webb, N.~M. (1991).
\newblock {\em Generalizability Theory: A Primer}.
\newblock Sage.

\bibitem[Vania et~al., 2021]{vania2021comparing}
Vania, C., et~al. (2021).
\newblock Comparing test sets with item response theory.
\newblock In {\em Proceedings of ACL}.

\bibitem[Yuan et~al., 2025]{yuan2025precisionpitfalls}
Yuan, Z., et~al. (2025).
\newblock Precision pitfalls in LLM evaluation.
\newblock {\em arXiv preprint}.

\end{thebibliography}

\end{document}
